\documentclass[12pt, letterpaper]{article} 

\usepackage{amsmath, amssymb, amsthm, enumerate}

\title{MATH 114 - Computational Stochastics}
\author{Dani Nguyen}
\date{Winter 2025}

\begin{document}
\maketitle
\tableofcontents

\section{Overview}
\subsection{Introduction}
Stochastic models describe real life phenomena like motion of the cell, shape of polymers
stock prices, and more. What can we compute \underline{analytically} and \underline{numerically} from these models? \\

\noindent Key questions:
\begin{enumerate}
    \item What is the distribution of the system?
    \item How can we sample from the distribution?
    \item How can we compute mean, variance, and higher moments of these systems?
\end{enumerate} 
\noindent Main topics:
\begin{enumerate}
    \item Sampling RVs \\
    \item Markov chain Monte Carlo (MCMC)
    \item Stochastic differential equations (numerical \& computation)
\end{enumerate}

\subsection{Sampling}
- Generating "hard" RVs from "easy" RVs. \\
- Compute "hard" integrals. \\ \\
\boxed{ex} Area of the unit circle, $A = \pi$ \\ 
Let $U_1 \sim U(-1,1)$ and $U_2 \sim U(-1,1)$ \\
Let $x = \begin{cases}
    0 \text{ if } (u_1,u_2) \text{ outside circle} \\
    1 \text{ if } (u_1,u_2) \text{ inside circle}
\end{cases}$ \\ 
$\implies  P(X = 1) = \frac{A}{4}, P(X = 0) = 1 - \frac{A}{4}$ \\ 
$\mathbb{E}(X) = \frac{A}{4}$  \\ 
After N trials we generate $X_1, X_2, ..., X_n$\quad iid \\
By the strong law of large numbers
\begin{align*}
    \lim_{N\rightarrow\inf} \frac{1}{N} \sum_{i=1}^{N}X_i = \frac{A}{4} = \frac{\pi}{4}
\end{align*}

\subsection{MCMC}
Construct Markov chains to generate target distribution. \\ 
\boxed{ex} Ising model


\section{Markov Chains}
\underline{Def:} A stochastic process is a family of RVs, $X_\epsilon \in S$, 
indexed by $t$, where $S \subseteq \mathbb{R}^d$ is a state space. Typically, 
$t \in [0,\infty)$, or $t \ in [a,b]$, or $t\in\{0,1,2,3,...\}$, etc.
\begin{itemize}
    \item A discrete time stochastic process, is a sequence of RVs, $X_n \in S$,
        $n = 0,1,2,...$ or $n = 1,2,3,...$
    \item If S is finite or countably infinite then the stochastic process is a 
        discrete state process.
\end{itemize}
First, we will examine discrete-time, discrete pace, stochastic process. \\
Two basic questions:
\begin{itemize}
    \item Find the joint distribution
    \begin{align*}
        &P(X_0 = i_0, X_1 = i_1, ..., X_n = i_n) \\
        &i_0, i_1, ..., i_n \in S
    \end{align*}
    \item Find the long-time behavior of the distribution of $X_n$:
    \begin{align*}
        P(x_n = i), i\in S, n>>1
    \end{align*}
\end{itemize}

\noindent \underline{Def:} Markov processes
\begin{itemize}
    \item A stochastic process $\{X_n\}^\infty_{n=0}$ is a Markov process if 
    \begin{align*}
        &P(X_{n + 1} = i_{n+1} | x_0 = i_0, ..., x_n = i_n) \\
        &= P(X_{n+1} i_{n+1} | X_n = i_n) \\
        & \forall n\ge 0, \forall i_0, ..., i_{n+1} \in S
    \end{align*}
    \item A Markov process is time homogeneous if: 
    \begin{align*}
        &P(x_{n+1} = j | x_n = i) \\ 
        &= P(X_1 = j | X_0 = i) \\
        &\forall n\ge 0, \forall i,j \in S
    \end{align*}
    Here we consider discrete time/state, time homogeneous Markov Chains (MC).
\end{itemize}

\noindent \underline{Def:} Transition probability
\begin{align*}
    P(i,j) &= P(X_1 = j | X_0 = i)\\
    &= P(i\rightarrow j ) \quad \forall i,j \in S
\end{align*}

\noindent \underline{Def:} Transition matrix
\begin{align*}
    &P=[p(i,j)] \text{ if } S = \{1,2,3,...\} \text{ or } \{0,1,2,...\} \\
    &\text{Here i,j are states and also are row/column labels}
\end{align*}

\noindent For any $n\ge 1, i_0, ..., i_n \in S$, we have:
\begin{align*}
    &P(X_0 = i_0, ..., X_n = i_n) \\
    &= P(X_0 = i_0) P(X_1 = i_1 | X_0 = i_0) P(X_2 = i_2 | X_1 = i_1, X_0 = i_0) \\
    &\ ...\ P(X_n = i_n | X_{n-1} = i_{n-1},..., X_0 = i_0)\qquad [\text{conditional probability}]
\end{align*}

\end{document}
